\documentclass[11pt]{article}
\usepackage[T1]{fontenc}
\usepackage[utf8]{inputenc}
\usepackage{lmodern}
\usepackage[margin=1in]{geometry}
\usepackage{amsmath,amssymb,amsthm,mathtools}
\usepackage{booktabs,longtable,array}
\usepackage{microtype}
\usepackage{xurl}
\usepackage[hidelinks]{hyperref}
\usepackage{enumitem,fancyvrb}
\allowdisplaybreaks[2]
\setlength{\emergencystretch}{2em}
\setcounter{tocdepth}{1}
\setlist[enumerate]{itemsep=3pt,topsep=5pt}
\newtheorem{theorem}{Theorem}[section]
\newtheorem{proposition}[theorem]{Proposition}
\newtheorem{lemma}[theorem]{Lemma}
\newtheorem{corollary}[theorem]{Corollary}
\theoremstyle{definition}
\newtheorem{definition}[theorem]{Definition}
\theoremstyle{remark}
\newtheorem{remark}[theorem]{Remark}
\numberwithin{equation}{section}
\newcommand{\R}{\mathbb R}
\newcommand{\Q}{\mathbb Q}
\newcommand{\C}{\mathbb C}
\newcommand{\Z}{\mathbb Z}
\newcommand{\N}{\mathbb N}
\newcommand{\abs}[1]{\left|#1\right|}
\newcommand{\norm}[1]{\left\|#1\right\|}
\newcommand{\PW}{\mathcal P_W}
\newcommand{\JV}{\mathcal J_V}
\newcommand{\OO}{\mathcal O}
\newcommand{\ELB}{\operatorname{ELB}}
\newcommand{\leanfile}[1]{\mbox{\texttt{#1}}}
\DeclareMathOperator{\ord}{ord}
\DeclareMathOperator{\lcm}{lcm}
\DeclareMathOperator{\codim}{codim}
\DeclareMathOperator{\Spec}{Spec}
\DeclareMathOperator{\Proj}{Proj}
\newcommand{\PaperAuthor}{Rohit Kumar Jha}
% No institutional affiliation.

\title{Optimal irrationality measures for logarithms of rational numbers\\
and a compatible-embedding extension}
\author{\PaperAuthor}
\date{October 8, 2026}
\hypersetup{
 pdftitle={Optimal irrationality measures for logarithms of rational numbers and a compatible-embedding extension},
 pdfauthor={\PaperAuthor},
 pdfsubject={Diophantine approximation and formal verification},
 pdfkeywords={logarithms, irrationality exponent, number fields, interpolation determinants, Lean}
}
\begin{document}
\maketitle
\begin{abstract}
We prove that the irrationality exponent of \(\log r\) is \(2\) for every
positive rational number \(r\ne1\), in particular for \(r=3\).
The proof extends the separated-weight interpolation and determinant
method of OpenAI's proof that the irrationality exponent of \(\pi\) is two.
The geometric extension permits distinct nonzero multiplicative centers,
and the arithmetic extension accounts for denominators of rational and
algebraic bases. More generally, let \(F\) be a number field of degree
\(d\), let \(\gamma,\xi\in F\) with \(\gamma\ne0\), and let
\(x\in\R\setminus\{0\}\). If
\(\exp(\sigma(\gamma)x)=\sigma(\xi)\) for \(s\ge1\) embeddings
\(\sigma:F\hookrightarrow\C\), then \(x\) is irrational and
\(2\le\mu(x)\le2d/s\).
The same algebraic determinant is estimated at every embedding; the field
norm supplies the factor \(d/s\). The root-of-unity case is included.
Consequences include a bound \(2\deg_{\Q}(\alpha)\) for
\(\mu(\log\alpha)\) when \(\alpha>0\) is algebraic and not \(1\),
and exponent two for rational-power logarithms, nonzero rational
arctangents, and certain quadratic-normalized logarithms.
The endpoint statements and their full dependency closure have been
checked in Lean, including replay in an initially empty kernel
environment. The proof does not compute numerical denominator thresholds.
\end{abstract}
\noindent\textit{Keywords:} irrationality exponent; logarithms; weighted
interpolation; algebraic number fields; formal verification.

\tableofcontents
\section{Introduction}

For an irrational real number \(x\), write
\[
 \mu(x)=\sup\left\{\nu>0:
  0<\abs{x-\frac pq}<q^{-\nu}
  \text{ for infinitely many reduced }p/q,\ q\ge2\right\},
\]
allowing the value \(+\infty\). Dirichlet's theorem gives
\(\mu(x)\ge2\). An upper bound concerns every sufficiently large
denominator, not merely a chosen sequence of approximants.
All logarithms of positive real numbers in this paper are natural.

\begin{theorem}[Rational logarithms]\label{thm:rational}
Let \(r\in\Q\), \(r>0\), \(r\ne1\). For every \(\varepsilon>0\)
there is an integer \(Q=Q(r,\varepsilon)\ge2\) such that
\begin{equation}\label{eq:rational-main}
 \abs{\log r-\frac pq}\ge q^{-(2+\varepsilon)}
 \qquad(p\in\Z,\ q\in\N,\ q\ge Q).
\end{equation}
No coprimality condition is imposed. Consequently \(\log r\) is
irrational and
\[
 \mu(\log r)=2.
\]
In particular, \(\mu(\log3)=2\).
\end{theorem}

The broader statement records exactly how the number field enters.
Embeddings, including two complex-conjugate embeddings, are counted
separately.

\begin{theorem}[Compatible embeddings]\label{thm:compatible}
Let \(F\) be a number field of degree \(d=[F:\Q]\), let
\(\gamma,\xi\in F\), \(\gamma\ne0\), and let \(x\in\R\setminus\{0\}\).
Suppose that a nonempty set \(S\) of embeddings \(F\hookrightarrow\C\)
satisfies
\begin{equation}\label{eq:compatibility}
 \exp(\sigma(\gamma)x)=\sigma(\xi)\qquad(\sigma\in S).
\end{equation}
Put \(s=|S|\) and \(\kappa=d/s\).
For every \(\varepsilon>0\) there is an integer \(Q\ge2\) such that
\begin{equation}\label{eq:compatible-main}
 \abs{x-\frac pq}\ge q^{-(2\kappa+\varepsilon)}
 \qquad(p\in\Z,\ q\in\N,\ q\ge Q).
\end{equation}
In particular \(x\) is irrational and
\[
 2\le\mu(x)\le\frac{2d}{s}.
\]
If all embeddings are compatible, then \(\mu(x)=2\).
\end{theorem}

Compatibility already implies \(\xi\ne0\). There is no hypothesis that
\(\xi\) is not a root of unity, that \(x\) is already irrational, or
that an interpolation matrix has full rank. Those facts or alternatives
needed in the proof are established below.

\begin{corollary}[Positive algebraic bases]\label{cor:algebraic}
If \(\alpha>0\) is algebraic over \(\Q\), \(\alpha\ne1\), and
\(d=\deg_{\Q}(\alpha)\), then \(\log\alpha\) is irrational and
\[
 2\le\mu(\log\alpha)\le2d.
\]
For every \(\varepsilon>0\), the associated lower bound with exponent
\(2d+\varepsilon\) holds for all sufficiently large denominators.
\end{corollary}
\begin{proof}
Take \(F=\Q(\alpha)\subset\R\), \(\gamma=1\),
\(\xi=\alpha\), \(x=\log\alpha\), and let \(S\) contain its given real
embedding into \(\C\). Theorem~\ref{thm:rational} is the case
\(F=\Q\), \(s=d=1\).
\end{proof}

\subsection{Prior work and the contribution here}

Quantitative approximation to logarithms has a substantial history.
Zudilin's survey \cite{Zudilin2004} discusses integral and hypergeometric
methods for \(\log2\), \(\log3\), and \(\pi\).
For \(\log3\), Wu and Wang \cite[Corollary~1]{WuWang2014} obtained
\(\mu(\log3)\le5.1163051\), subsequently improved by Bondareva,
Luchin, and Salikhov \cite[Corollary~1]{BLS2018} to
\(\mu(\log3)\le5.116201\). These are finite upper bounds,
not the exact value two asserted in Theorem~\ref{thm:rational}.
Marcovecchio and Viola \cite{MV2012} study approximation to logarithms
by elements of number fields, as well as nonquadraticity measures,
using double integrals and the saddle-point method.
Their number-field approximation problem, normalized by Weil height,
must be distinguished from the rational-approximation exponent used here.
Marcovecchio's multiple-Legendre construction \cite{Marcovecchio2014}
also treats approximation by algebraic numbers of bounded degree.
Bouchelaghem, He, Li, and Wu \cite{BHLW2024} study linear-independence
measures involving two rational logarithms, another quantitative problem
distinct from the individual exponents considered here.
We do not claim that finite irrationality bounds for logarithms were
previously unknown, nor that our degree bound is best in every special
case.
The immediate input is OpenAI's preprint \cite{OAIpi} and its pinned
formal development. Its proof of
\(\mu(\pi)=2\) combines separated approximation denominators, weighted
interpolation, arithmetic clearing of a nonzero determinant, and an
analytic saving caused by repeated Taylor coefficient rows.
This paper reuses that architecture and its proved general ingredients.
It is not a deduction from the numerical statement \(\mu(\pi)=2\).

The new geometric input allows the multiplicative coordinates of the
interpolation points to be distinct nonzero numbers. For \(\log3\)
they are \(3^j\); for \(\log r\), they are \(r^j\).
The local coordinate \(t=Y/y_j-1\) preserves the logarithmic frame.
Once the rigidity argument forces \(Y\) to be constant, such a curve
can meet at most one center. The formal development also supplies the
common-fiber alternative needed when a root of unity is reduced to \(1\).
Both alternatives lead to the actual, explicitly defined interpolation
matrix, not to a matrix of assumed rank.

The arithmetic extension clears denominators of both the algebraic base
\(\xi\) and the slope \(\gamma\). A single selected determinant over
\(F\) is then mapped through every embedding. Compatible embeddings
give analytic decay; the others incur a controlled growth cost.
Taking its norm and dividing by the number of compatible embeddings
produces \(\kappa=d/s\). Omitting the other embeddings would erase this
essential factor and would not prove the stated theorem.

\subsection{Organization and scope}

Section~\ref{sec:preliminaries} fixes the approximation conventions.
Section~\ref{sec:interpolation} proves the required interpolation
statement, making the inherited multiplicity input explicit.
Sections~\ref{sec:arithmetic} and \ref{sec:analysis} establish the
arithmetic and analytic inequalities for the same determinant.
Section~\ref{sec:parameters} chooses all parameters in their required
order and closes the contradiction.
Section~\ref{sec:consequences} gives the specializations, a comparison
with the \(\pi\) proof, and counterexamples to unrestricted real-input
claims. Section~\ref{sec:verification} records formal statements,
source provenance, and reproducible checks.

The results do not assert exponent two for every positive algebraic
base or a uniform finite bound for arbitrary positive real bases.
The eventual thresholds are existential; this paper does not compute
their numerical values.

\section{Approximation conventions and preliminary reductions}
\label{sec:preliminaries}

\begin{definition}
For \(x,b\in\R\), write \(\ELB(x,b)\) if, for every real \(\nu>b\),
there is \(Q\in\N\), \(Q\ge2\), such that
\[
 \forall p\in\Z\ \forall q\in\N,\qquad
 q\ge Q\ \Longrightarrow\ q^{-\nu}\le\abs{x-p/q}.
\]
\end{definition}

\begin{lemma}\label{lem:endpoint}
If \(\ELB(x,b)\), then \(x\) is irrational. Moreover, for
\(b\ge2\),
\[
 2\le\mu(x)\le b.
\]
\end{lemma}
\begin{proof}
If \(x=a/c\), with \(a\in\Z\), \(c\in\N_{>0}\), use
\(p=ka\), \(q=kc\) with \(k\) arbitrarily large. This contradicts
the positive lower bound in the definition.
For any \(\nu>b\), that definition excludes good reduced approximants
with \(q\ge Q\). There are only finitely many smaller denominators
and finitely many relevant numerators for each of them. Thus no
approximation exponent greater than \(b\) belongs to the defining set
of \(\mu(x)\). Dirichlet supplies the opposite inequality.
\end{proof}

It is enough to contradict the following assertion for each fixed
\(\nu>2\kappa\):
\begin{equation}\label{eq:bad-approximations}
 \text{for every }Q,\quad
 \exists p\in\Z,\ q\in\N,\quad
 q\ge Q,\qquad\abs{x-p/q}\le q^{-\nu}.
\end{equation}
Failure of \(\ELB(x,2\kappa)\) implies this assertion for some such
\(\nu\). Since \(x\ne0\), sufficiently large approximants in
\eqref{eq:bad-approximations} have \(p\ne0\).
They also satisfy \(\abs{p/q}\le |x|+1\).
The integers
\[
 w_i=\lceil\log q_i\rceil
\]
can be made successively as large as desired. In particular, one may
require each \(w_i\) to exceed a threshold depending on the previously
chosen weights. This does not require a prescribed growth rate of the
whole sequence of good approximants.

\begin{lemma}[Root-of-unity reduction]\label{lem:torsion}
In proving Theorem~\ref{thm:compatible}, one may assume either that
the powers \(\xi^j\), \(j\ge0\), are distinct or that \(\xi=1\).
The reduction preserves \(x,F,S,d,s\).
\end{lemma}
\begin{proof}
If \(\xi\) is not a root of unity, equality of two distinct powers
would give a positive power equal to \(1\), since \(\xi\ne0\).
If \(\xi^n=1\) for some \(n\ge1\), replace \((\gamma,\xi)\)
by \((n\gamma,1)\). Equation~\eqref{eq:compatibility} remains true
after raising both sides to the \(n\)-th power.
The new slope is nonzero.
\end{proof}

From now through Section~\ref{sec:parameters}, this normalization has
been performed. Its effect is only on fixed constants.
All choices other than the final degree parameter \(H\) are fixed
before any limit \(H\to\infty\) is taken.
For multi-indices \(\alpha\in\N^m\), put
\[
 |\alpha|=\sum_{i=1}^m\alpha_i,\qquad
 w\alpha=\sum_{i=1}^m w_i\alpha_i,\qquad
 \binom{\alpha}{\beta}=\prod_{i=1}^m\binom{\alpha_i}{\beta_i}.
\]
Coordinatewise inequalities between multi-indices have their usual
meaning. Empty products and empty least common multiples equal \(1\).

\section{Interpolation at logarithmic centers}\label{sec:interpolation}

Let \(m,K\ge1\), and fix positive rational \(w_0,v_0,\theta\)
with \(0<\theta<1\). For further positive rational weights \(w_i\),
set
\[
 W=(w_0,w_1,\ldots,w_m),\qquad
 V=(v_0,v_1,\ldots,v_m),\qquad v_i=w_i/\theta\quad(i>0).
\]
Let \(\PW(H)\) be the space spanned by \(Y^hX^\alpha\) with
\(w_0h+w\alpha\le H\). Let \(\JV(H)\) be the quotient of
\(\C[[t,u_1,\ldots,u_m]]\) retaining precisely the coefficients
\[
 v_0\ell+\sum_i v_i\beta_i<H.
\]
The inequality for rows is strict, whereas the inequality for columns
is weak.

\begin{theorem}[Separated logarithmic interpolation]\label{thm:interpolation}
Assume
\begin{equation}\label{eq:volume}
 K(w_0/v_0)\theta^m<1.
\end{equation}
There are successive lower thresholds for \(w_1,\ldots,w_m\),
independent of the center coordinates, with this property.
For nonzero \(y_j\in\C\), \(0\le j<K\), and \(c_{ji}\in\C\),
the map
\begin{equation}\label{eq:jetmap}
 \begin{split}
 \PW(H)&\longrightarrow\bigoplus_{j<K}\JV(H),\\
 P&\longmapsto
 \left(P\bigl(y_j(1+t),
       (c_{ji}+u_i+\log(1+t))_{i=1}^m\bigr)\right)_{j<K}
 \end{split}
\end{equation}
is surjective for all sufficiently large \(H\) divisible by a fixed
positive integer, in either of the following cases:
\begin{enumerate}[label=(\roman*)]
\item the \(y_j\) are pairwise distinct;
\item all \(y_j=1\), each family \(c_{0i},\ldots,c_{K-1,i}\)
is pairwise distinct, and \(K\theta^m<1\).
\end{enumerate}
The divisibility requirement depends only on the weights. The lower
threshold for \(H\) may depend on every fixed datum, including the
centers.
\end{theorem}

Case (ii) is the common-fiber interpolation of \cite[Section 2]{OAIpi}.
We give the argument with both cases present to exhibit the exact
change required for case (i). In particular, a count of coefficients
alone is not used to assert surjectivity.

\subsection{The uniform local input and separated weights}

We use the weighted transverse multiplicity comparison of
\cite[Lemma 2.2]{OAIpi}. In the notation needed here it says the following.

\begin{lemma}[Weighted multiplicity comparison]\label{lem:multiplicity}
Give \(D\) affine coordinates positive degree weights \(\rho_a\).
Let commuting analytic vector fields form a frame on an open set,
with positive derivative costs \(\kappa_b\).
Suppose polynomials \(f_\ell\) have weighted degree at most \(N\),
and \(Z\) is a codimension-\(k\) component of their common zero set.
If every further frame derivative of cost less than
\(\epsilon N\) of every \(f_\ell\) vanishes on \(Z\), then, for
any coordinate normal basis \(A\) and frame normal basis \(B\),
\begin{equation}\label{eq:local-comparison}
 \prod_{a\in A}\rho_a
 \le 2k!\epsilon^{-k}\prod_{b\in B}\kappa_b.
\end{equation}
The constant is independent of the equations and of the component.
\end{lemma}
\begin{proof}
Work at a general smooth point of \(Z\) where both normal-basis
conditions hold and no other component passes through the point.
The commuting flows in the \(B\)-directions give transverse analytic
coordinates. Vanishing of the derivatives says that the equations
have no transverse terms of weight below \(\epsilon N\).
Restrict to the coordinate slice in the \(A\)-directions.
That slice is transverse, its common zero is isolated, and substitution
of the remaining analytic coordinates cannot introduce a transverse
term of smaller weight.
Thus the local quotient length is at least
\[
 \#\{\alpha\in\N^k:\textstyle\sum_{b\in B}\kappa_b\alpha_b
                  <\epsilon N\}
 \ge \frac{(\epsilon N)^k}{2k!\prod_{b\in B}\kappa_b}.
\]
For the last inequality, unit cubes indexed by the integer parts of
points of the open weighted simplex cover that simplex; even the
factor \(2\) is unnecessary for this count.

Generic \(k\) constant linear combinations of the restricted equations
still cut out an isolated zero. Their quotient length bounds the
preceding length from above, and weighted B\'ezout bounds it by
\(N^k/\prod_{a\in A}\rho_a\).
For clarity, this weighted version follows by choosing an integer \(M_0\)
with \(M_0\rho_a\in\N_{>0}\) and substituting
\(x_a=\lambda_a+z_a^{M_0\rho_a}\). Choose \(\lambda_a\) so that
the preimages of the point are unramified. There are
\(\prod_{a\in A}M_0\rho_a\) such preimages, each preserving the
local length, while the total degrees are at most \(M_0N\).
Ordinary B\'ezout for isolated zeros gives the asserted upper bound;
see \cite[Corollary VIII.3]{Mondal2021}.
Comparison of the two lengths proves \eqref{eq:local-comparison}.
\end{proof}

Choose \(\sigma_0\in\Q_{>0}\) small enough that
\begin{equation}\label{eq:sigma}
 \begin{split}
 (1+3\sigma_0)^{m+1}K(w_0/v_0)\theta^m&<1,\\
 (1+\sigma_0)\theta&<1,
 \end{split}
\end{equation}
and, in case (ii), also
\((1+3\sigma_0)^mK\theta^m<1\).
Put \(L_0=m+2\) and take the deliberately enlarged constant
\[
 C_*=
 2(m+2)^{m+2}\left(1+\frac{m+2}{\sigma_0}\right)^{m+2}.
\]
It dominates \(2k!(L_0/\sigma_0)^k\) for \(1\le k\le m\).
Choose the positive weights successively so that
\begin{equation}\label{eq:weight-separation}
 \prod_{a\in A}w_a>C_*\prod_{b\in B}v_b
\end{equation}
whenever \(A,B\subseteq\{0,\ldots,m\}\) have equal size between
\(1\) and \(m\), and the largest positive index in their symmetric
difference belongs to \(A\).
At that index the ratio in \eqref{eq:weight-separation} contains the
new weight linearly. Larger common indices contribute only
\(w_i/v_i=\theta\); all remaining factors were already fixed.
Finitely many inequalities therefore impose a finite lower threshold.
No point coordinate enters these thresholds.

The formal implementation can use a rectangular transverse-length
bound, with constant \(k^k(L_0/\sigma_0)^k\), at a sufficiently large
auxiliary degree. The displayed \(C_*\) also dominates that constant.
The auxiliary-degree threshold is independent of the varying component,
so either local formulation gives the same separated-weight argument.

\subsection{The curve inequality}

On the smooth complete model of an irreducible curve \(C\) meeting
the affine chart, define
\begin{equation}\label{eq:weighted-degree}
 \deg_W C=\sum_P
  \max\left\{0,-\frac{\ord_PY}{w_0},
          -\frac{\ord_PX_1}{w_1},\ldots,
          -\frac{\ord_PX_m}{w_m}\right\}.
\end{equation}
An identically zero coordinate contributes no pole.
At a branch through \(a_j=(y_j,c_{j1},\ldots,c_{jm})\), let
\begin{equation}\label{eq:contact}
 h_P=\min\left\{
  \frac{\ord_P(Y/y_j-1)}{v_0},
  \frac{\ord_P(X_i-c_{ji}-\log(Y/y_j))}{v_i}:1\le i\le m
 \right\}.
\end{equation}
The logarithm is a formal series at the branch, not a global choice
of a complex logarithm. These contacts are positive and finite on
a nonconstant branch.

\begin{proposition}\label{prop:curve}
For the separated weights above, in either case of
Theorem~\ref{thm:interpolation},
\begin{equation}\label{eq:curve-bound}
 \deg_W C\ge(1+\sigma_0)\sum_{P\mapsto a_j}h_P.
\end{equation}
\end{proposition}
\begin{proof}
Suppose the inequality fails. A count of coefficients and conditions
gives, for large auxiliary integers \(N\), a nonzero polynomial \(F_N\)
of \(W\)-degree at most \(N\), with local \(V\)-order at least
\((1+3\sigma_0)N\) at every center. The leading terms in the two counts
are
\[
 \frac{N^{m+1}}{(m+1)!\prod_aw_a},
 \qquad
 \frac{K(1+3\sigma_0)^{m+1}N^{m+1}}
      {(m+1)!\prod_av_a},
\]
and \eqref{eq:sigma} makes the first larger. Independence of the
conditions is not needed to produce \(F_N\).

On \(Y\ne0\), use the commuting frame
\[
 D_0=Y\partial_Y+\sum_{i=1}^m\partial_{X_i},
 \qquad D_i=\partial_{X_i}\quad(i>0).
\]
In the coordinates of \eqref{eq:jetmap}, these are
\((1+t)\partial_t,\partial_{u_i}\).
Each preserves the bound on polynomial degree; differentiating with
cost \(v_a\) lowers local weighted order by at most \(v_a\).
Thus every derivative of \(F_N\) of total cost at most
\(\sigma_0N\) has branch order at least
\((1+2\sigma_0)Nh_P\).
Its total pole order on the complete curve is at most \(N\deg_W C\).
The assumed failure of \eqref{eq:curve-bound} makes the former
total strictly larger. The zero--pole equality forces all these
derivatives to vanish identically on \(C\).

Set \(\delta_N=\sigma_0N/L_0\). Consider successively the common
zero sets of derivatives of cost at most \(r\delta_N\),
\(0\le r\le L_0\), on \(Y\ne0\).
Choose nested irreducible components containing the affine part of
\(C\), working backwards from the last level.
Their dimensions lie between \(1\) and \(m\); hence two adjacent
components coincide. Call this persistent component \(Z\), of
codimension \(k\), \(1\le k\le m\).
Further derivatives of cost less than \(\delta_N\) of its defining
equations from the earlier level vanish on \(Z\).
Lemma~\ref{lem:multiplicity} gives
\[
 \prod_{a\in A}w_a\le C_*\prod_{b\in B}v_b
\]
for every coordinate normal basis \(A\) and frame normal basis \(B\).
This argument is an instance of the persistent derivative-ideal method;
see also \cite[Section 5]{Philippon1986}.

If some positive coordinate direction is nontangent to \(Z\),
choose its largest index \(i\). Some coordinate normal basis contains
\(i\). Separation forces every frame normal basis to contain \(D_i\).
The span of the remaining frame vectors is
\(\ker(dX_i-dY/Y)\). Its failure to span the normal quotient implies
\[
 dX_i=dY/Y\quad\text{on }Z.
\]
This forces \(Y\) to be constant. Otherwise restrict to an algebraic
curve on which \(Y\) is nonconstant and pass to its complete model.
At a zero or pole of \(Y\), the residue of \(dY/Y\) is the
nonzero integer \(\ord_PY\), whereas the exact differential \(dX_i\)
has residue zero. This is impossible.
If every positive coordinate direction is tangent, properness of \(Z\)
forces its tangent space to be their span, so \(dY=0\) and the same
conclusion holds. In particular \(Y\) is constant on \(C\).

In case (i), \(C\) meets at most one prescribed center \(a_j\).
Choose \(X_\ell\) nonconstant on \(C\). The function
\(X_\ell-c_{j\ell}\) has total pole degree at most
\(w_\ell\deg_W C\), and has zero order at least \(v_\ell h_P\)
at every relevant branch. Consequently
\[
 \deg_W C\ge\theta^{-1}\sum_{P\mapsto a_j}h_P
          >(1+\sigma_0)\sum_{P\mapsto a_j}h_P,
\]
contradicting the supposition.

In case (ii), the constant is \(Y=1\).
In that fiber a second coefficient count is available because
\(K(1+3\sigma_0)^m\theta^m<1\).
It gives a nonzero polynomial in \(X_1,\ldots,X_m\) of degree
at most \(N\) and order at least \((1+3\sigma_0)N\) at all centers.
The same zero--pole argument makes its low-cost ordinary partial
derivatives vanish on \(C\).
For \(m=1\), a nonzero polynomial cannot vanish on this curve.
For \(m\ge2\), repeat the persistent-component argument in the fiber.
The coordinate and derivative frames now agree.
Two different normal-basis index sets would violate
\eqref{eq:weight-separation}, so there is just one.
Choose an index \(i\) in it. The other coordinate directions fail
to span the normal quotient, so \(dX_i=0\) on the component.
That coordinate is constant in characteristic zero.
Coordinatewise distinctness implies that the component, and hence
\(C\), meets at most one center. The preceding single-center
zero--pole argument again contradicts the supposed excess.
\end{proof}

\subsection{From curve degrees to the actual coefficient map}

We finish the proof of Theorem~\ref{thm:interpolation}; this is the
curve-to-jet passage of \cite[Section 2.3]{OAIpi}, with normalized
coordinates at arbitrary \(y_j\ne0\).
Choose an integer \(R\) for which every \(R/w_a\) and \(R/v_a\)
is a positive integer. Embed the affine space by all monomials of
\(W\)-degree at most \(R\), including \(1\), and take its projective
closure \(X\), with \(L=\OO_X(1)\).
Pure coordinate powers attain the maximum pole at each point of a
complete curve, so
\begin{equation}\label{eq:projective-degree}
 \deg(L|_C)=R\deg_W C.
\end{equation}

Replace each logarithm by a Taylor polynomial \(G_i(t)\) whose first
omitted term has \(t\)-weight greater than \(v_i\). The coordinates
\[
 t=Y/y_j-1,\qquad u_i'=X_i-c_{ji}-G_i(t)
\]
are algebraic regular coordinates at \(a_j\).
The replacement preserves \eqref{eq:contact} and induces inverse
filtration-preserving triangular changes on every finite coefficient
quotient. Define the finite-colength ideal
\[
 I_j=(t^{R/v_0},(u_1')^{R/v_1},\ldots,(u_m')^{R/v_m}).
\]
Patch these ideals with the unit ideal away from the centers.
Its branch order is \(Rh_P\).

Let \(p:X'\to X\) be the ordinary blow-up of this ideal \(I\), and
write \(I\OO_{X'}=\OO_{X'}(-E)\), \(A=p^*L\).
The curve inequality and \eqref{eq:projective-degree} show that
\(A-(1+\sigma_0)E\) is nef.
For a curve contracted by \(p\), this follows instead from relative
ampleness of \(-E\); curves avoiding the affine chart do not meet \(E\).
For large \(a>1\), \(aA-E\) is ample.
The identity
\[
 A-E=
 \frac{a-1}{a(1+\sigma_0)-1}\bigl(A-(1+\sigma_0)E\bigr)
 +\frac{\sigma_0}{a(1+\sigma_0)-1}(aA-E)
\]
therefore makes \(A-E\) ample; see
\cite[Corollary 1.4.10]{Lazarsfeld2004}.
These statements do not require a smooth or normal blow-up.

The large-degree graded-piece and vanishing theorems for the Rees
algebra give
\[
 p_*\OO_{X'}(-nE)=I^n,\qquad
 R^q p_*\OO_{X'}(-nE)=0\quad(q>0,\ n\text{ sufficiently large}).
\]
These are ordinary powers, not integral closures.
The relevant general facts are \cite[Tags 0AG6 and 0AG7]{Stacks}.
Projection, Leray, and Serre vanishing now give
\[
 H^1(X,I^n\otimes L^n)=
 H^1(X',\OO_{X'}(n(A-E)))=0
\]
for large \(n\).
Hence global sections of \(L^n\) realize every local class modulo
\(I^n\). Each generator of \(I\) has weight \(R\), so
\[
 I^n\subseteq\{\text{series of weight at least }nR\}.
\]
Surjectivity onto the former quotient therefore implies surjectivity
onto the desired smaller coefficient quotient.
Completion does not change the finite-length local quotients.
Finally, for sufficiently large \(n\), global sections of \(L^n\)
are restrictions of homogeneous polynomials of degree \(n\) in the
embedding coordinates, by Serre vanishing for the projective ideal
sheaf. They have \(W\)-degree at most \(nR\) on the affine chart.
Taking \(H=nR\) and undoing the triangular logarithmic substitution
proves \eqref{eq:jetmap}.
\hfill\(\square\)

\section{The algebraic determinant and its arithmetic bound}
\label{sec:arithmetic}

Fix \(\nu>2\kappa\) and assume \eqref{eq:bad-approximations}.
For now take finite parameters \(m,K,w_0,v_0,\theta\) satisfying the
interpolation hypotheses, and choose good approximants
\(p_i/q_i\), \(q_i\ge2\), \(p_i\ne0\), with separated
\(w_i=\lceil\log q_i\rceil\). Their precise order of choice is
given in Section~\ref{sec:parameters}.
Let \(F_0>2/\theta\), and set
\begin{equation}\label{eq:truncations}
 T_i=\left\lceil\frac{F_0w_i}{v_0}\right\rceil,\qquad
 G_i(t)=\sum_{1\le k<T_i}\frac{(-1)^{k+1}}{k}t^k,\qquad
 L_i=\lcm(1,\ldots,T_i).
\end{equation}
Using \(T_i\), rather than \(T_i-1\), in the least common multiple
is harmless and agrees with the formal arithmetic bound.
The omitted tail has weight at least \(v_0T_i\ge F_0w_i>w_i/\theta\).

\subsection{Matrix, rank, and one minor over the number field}

Index columns by \((h,\alpha)\in\N\times\N^m\) satisfying
\(w_0h+w\alpha\le H\), and rows by
\((j,\ell,\beta)\) satisfying
\[
 0\le j<K,\qquad v_0\ell+w\beta/\theta<H.
\]
Define the matrix over \(F\) by
\begin{equation}\label{eq:matrix}
 M_H[(j,\ell,\beta),(h,\alpha)]
 =[t^\ell u^\beta]\,
   \bigl(\xi^j(1+t)\bigr)^h
   \prod_{i=1}^m
    \bigl(j\gamma p_i/q_i+G_i(t)+u_i\bigr)^{\alpha_i}.
\end{equation}
Explicitly, for \(\beta\le\alpha\), its entry is
\begin{equation}\label{eq:binomial-entry}
 \xi^{jh}\binom{\alpha}{\beta}
 [t^\ell](1+t)^h
   \prod_i\bigl(j\gamma p_i/q_i+G_i(t)\bigr)^{\alpha_i-\beta_i};
\end{equation}
it is zero otherwise.

Fix one embedding in \(S\). If the powers of \(\xi\) are distinct,
the centers
\[
 \bigl(\sigma(\xi)^j,
       (j\sigma(\gamma)p_i/q_i)_{i=1}^m\bigr)
\]
satisfy case (i) of Theorem~\ref{thm:interpolation}.
If \(\xi=1\), each additive coordinate is injective in \(j\)
because \(\gamma p_i/q_i\ne0\), and case (ii) applies.
The truncation in \eqref{eq:truncations} preserves the jet quotient.
It follows that, for sufficiently large divisible \(H\), the
complex image of \eqref{eq:matrix} has full row rank.
Select a square minor using every row, and denote its determinant
by \(\Delta_H\in F^\times\) and its size by \(M=M(H)\).
The chosen columns may depend on \(H\).
Every embedding sends this same determinant to the determinant of
the corresponding complex minor; none of these images is zero.
We never choose different minors at different embeddings.

The fixed-weight simplex counts give
\begin{equation}\label{eq:row-count}
 M(H)\sim
 \frac{K\theta^m H^{m+1}}{(m+1)!v_0\prod_iw_i}.
\end{equation}
Define
\begin{equation}\label{eq:mean}
 \bar b_H=\frac1{MH}\sum_{(j,\ell,\beta)\ {\rm row}}w\beta,
 \qquad 0\le\bar b_H\le\theta.
\end{equation}
The simplex asymptotic follows, for example, by comparing unit cubes
with the bounded shifts of the corresponding real simplex. Changing
a weak boundary inequality to a strict one does not change its leading
volume.

\subsection{Clearing the base, slope, and logarithmic denominators}

Choose positive integers \(D_\xi,D_\gamma\) such that
\(D_\xi\xi,D_\gamma\gamma\in\OO_F\), and put
\begin{equation}\label{eq:common-denominator}
 \mathcal D_H=
 D_\xi^{\,K\lfloor H/w_0\rfloor}
 \prod_{i=1}^m(D_\gamma L_i)^{\lfloor H/w_i\rfloor}.
\end{equation}
Scale each column by \(\prod_iq_i^{\alpha_i}\), each row by
\(\prod_iq_i^{-\beta_i}\), and every entry by \(\mathcal D_H\).

\begin{lemma}\label{lem:integral}
Every scaled entry belongs to \(\OO_F\). Hence the scaled determinant
is a nonzero algebraic integer
\[
 Z_H=S_H\Delta_H,\qquad
 S_H=\mathcal D_H^M
   \prod_{\text{selected columns}}\prod_iq_i^{\alpha_i}
   \prod_{\text{rows}}\prod_iq_i^{-\beta_i}>0.
\]
Here \(S_H\in\Q\), and it need not be an integer.
\end{lemma}
\begin{proof}
Entries with \(\beta\not\le\alpha\) vanish.
For the others, the row and column factors turn
\eqref{eq:binomial-entry} into
\[
 \xi^{jh}\binom{\alpha}{\beta}
 [t^\ell](1+t)^h
   \prod_i(j\gamma p_i+q_iG_i(t))^{\alpha_i-\beta_i}.
\]
The polynomial
\(D_\gamma L_i(j\gamma p_i+q_iG_i(t))\)
has integral coefficients. Its required power is at most
\(\alpha_i\le\lfloor H/w_i\rfloor\).
Also \(jh\le K\lfloor H/w_0\rfloor\), so
\[
 D_\xi^{K\lfloor H/w_0\rfloor}\xi^{jh}
 =(D_\xi\xi)^{jh}D_\xi^{K\lfloor H/w_0\rfloor-jh}
\]
is integral. This proves entrywise integrality.
The determinant of an integral matrix is integral, and the scalings
are nonzero.
\end{proof}

Let \(\Lambda=\log4+4\). We use the elementary bound
\(\log\lcm(1,\ldots,n)\le\Lambda n\).
Indeed, the elementary Chebyshev bound
\(\vartheta(u)\le u\log4\), obtained from binomial coefficients,
and the prime-power identity for \(\psi\) give, for \(n\ge1\),
\[
 \log\lcm(1,\ldots,n)=\psi(n)
 \le n\log4+2\sqrt n\log n
 \le(\log4+4)n.
\]
The last step uses \(\log n\le2\sqrt n\); \(n=0\) is immediate.
This deliberately loose constant is exactly
\path{OAI.PiExponent.Arithmetic.lcmConstant} in the pinned source.
The sharper constant in the informal argument of
\cite[Section 3.1]{OAIpi} is not needed.
Since \(T_i\le F_0w_i/v_0+1\),
\begin{equation}\label{eq:log-D}
 \frac{\log\mathcal D_H}{H}
 \le \frac{\Lambda F_0m}{v_0}
    +(\Lambda+\log D_\gamma)\sum_i\frac1{w_i}
    +\frac{K\log D_\xi}{w_0}.
\end{equation}
Write \(w_*=\min_iw_i\). For every selected column,
\(\sum_i\alpha_i\log q_i\le H\).
For a row, \(\log q_i\ge w_i-1\) and
\(|\beta|\le\theta H/w_*\), so
\[
 \sum_i\beta_i\log q_i\ge w\beta-\theta H/w_*.
\]
Consequently
\begin{equation}\label{eq:scaling-bound}
 \frac{\log S_H}{MH}\le1-\bar b_H+E_{\rm ar}^{\,0},
 \qquad
 E_{\rm ar}^{\,0}=
 \frac{\Lambda F_0m}{v_0}
 +(\Lambda+\log D_\gamma)\sum_i\frac1{w_i}
 +\frac{K\log D_\xi}{w_0}+\frac{\theta}{w_*}.
\end{equation}

\begin{proposition}[Norm lower bound]\label{prop:arithmetic}
Set
\[
 \mathcal N_H=\frac1{sMH}
        \sum_{\sigma:F\hookrightarrow\C}\log|\sigma(\Delta_H)|,
 \qquad E_{\rm ar}=\kappa E_{\rm ar}^{\,0}.
\]
Then
\begin{equation}\label{eq:arithmetic}
 \mathcal N_H\ge-\kappa(1-\bar b_H)-E_{\rm ar}.
\end{equation}
\end{proposition}
\begin{proof}
The norm of the nonzero algebraic integer \(Z_H\) is a nonzero
rational integer. Therefore
\[
 1\le|N_{F/\Q}(Z_H)|
   =S_H^d\prod_{\sigma:F\hookrightarrow\C}|\sigma(\Delta_H)|.
\]
Take logarithms, divide by \(sMH\), and apply
\eqref{eq:scaling-bound}.
\end{proof}

For \(x=\log3\), take \(F=\Q\), \(\gamma=1\), \(\xi=3\),
and \(D_\xi=D_\gamma=1\): the factor \(3^{jh}\) is integral
and costs no denominator. For a positive rational \(r=a/b\),
one can take \(D_\xi=b\); the extra cost is at most
\(K\log b/w_0\), which the parameter choice will make small.

\section{Analytic bounds at all embeddings}\label{sec:analysis}

Define fixed size constants
\[
 \Gamma=\max\left(1,\max_\sigma|\sigma(\gamma)|\right),\quad
 \Xi=\max\left(1,\max_\sigma|\sigma(\xi)|\right),\quad
 B_0=\Gamma(|x|+1),\quad R_0=100(1+\Gamma|x|).
\]
The maxima range over all embeddings, not just \(S\).
Thus \(B_0\ge1\), \(R_0\ge100\). Put \(R=R_0K\).

\subsection{An exact row translation}

For a polynomial column \(P\) and \(a\in\N^m\), consider the
entire function
\begin{equation}\label{eq:entire}
 f_{a,P}(z)=[u^a]P(e^z,z+u_1,\ldots,z+u_m).
\end{equation}
For \(P=Y^hX^\alpha\), this is
\(\binom{\alpha}{a}e^{hz}z^{|\alpha|-|a|}\) when \(a\le\alpha\),
and zero otherwise. In particular these functions do not depend on
the center index \(j\).

Fix \(\sigma\in S\) and put
\[
 \lambda_\sigma=\sigma(\gamma)x,\quad
 z=j\lambda_\sigma+\log(1+t),\quad
 \epsilon_i=j\sigma(\gamma)(p_i/q_i-x),\quad
 \tau_i(t)=G_i(t)-\log(1+t).
\]
The two identities
\[
 e^z=\sigma(\xi)^j(1+t),\qquad
 j\sigma(\gamma)p_i/q_i+G_i(t)+u_i
       =z+u_i+\epsilon_i+\tau_i(t)
\]
give the exact row-vector expansion
\begin{equation}\label{eq:row-expansion}
 \begin{split}
 {\rm row}_{j,\ell,\beta}
 =\sum_{\substack{a\ge\beta\\wa\le H}}\binom a\beta
   \sum_{0\le e\le a-\beta}\binom{a-\beta}{e}
       \epsilon^{a-\beta-e}
   \sum_{k=0}^{\ell}
    \left([t^k]\prod_i\tau_i(t)^{e_i}\right)
    \left([t^{\ell-k}]
       f_{a,P}(j\lambda_\sigma+\log(1+t))\right)_P .
 \end{split}
\end{equation}
The scalar multiplying a translated row is independent of the column
\(P\). This is what permits determinant multilinearity.
The identity is an equality of finite coefficient calculations:
the column space bounds \(a\), and only Taylor coefficients through
\(\ell\) occur.

\begin{lemma}[Scalar cost]\label{lem:scalar}
Every scalar in \eqref{eq:row-expansion} has absolute value at most
\begin{equation}\label{eq:scalar}
 \exp\{-\nu w(a-\beta)+H E_{\rm tr}\},
 \qquad
 E_{\rm tr}=\frac{\nu}{F_0}+\frac{\log2}{v_0}
       +\frac{\log4+\log(2K\Gamma)+\nu}{w_*}.
\end{equation}
Each row has at most
\[
 Q_H=(\lfloor H\rfloor+1)^{2m}(\lfloor H/v_0\rfloor+1)
\]
choices in this expansion.
\end{lemma}
\begin{proof}
The approximation and ceiling inequalities give
\[
 |\epsilon_i|\le K\Gamma q_i^{-\nu}
              \le K\Gamma e^\nu e^{-\nu w_i}.
\]
Each \(\tau_i\) starts at degree \(T_i\).
A nonzero coefficient in its product requires
\(\sum_i e_iT_i\le k\le\ell<H/v_0\), whence
\(we<H/F_0\).
On \(|t|=1/2\), the logarithmic tails have absolute value at most
\(\sum_{n\ge1}2^{-n}/n=\log2<1\).
Cauchy's estimate bounds the coefficient of their product by \(2^k\).
The binomial coefficients contribute at most \(4^{|a|}\), and
\(|a|\le H/w_*\).
In the approximation factor, replacing \(w(a-\beta-e)\) by
\(w(a-\beta)\) costs at most \(\nu H/F_0\).
These inequalities imply \eqref{eq:scalar}, with the harmless extra
factor \(2\) inside \(\log(2K\Gamma)\).
Finally \(w_i\ge1\), so both \(a\) and \(e\) have coordinates at
most \(\lfloor H\rfloor\); \(k\le\lfloor H/v_0\rfloor\).
\end{proof}

\subsection{The repeated-coefficient saving}

Select one translated row from each original row and omit its scalar.
Call the resulting square matrix \(T\). Let \(n_a\) count the chosen
rows with transverse index \(a\), so \(\sum_a n_a=M\).
Put
\[
 c=\frac{\log2}{4},\qquad
 E_{\rm hol}=\frac{R_0K}{w_0}+\frac{\log2}{v_0}
                   +\frac{\log(2R_0K)}{w_*}.
\]

\begin{lemma}[Collision estimate]\label{lem:collision}
Uniformly over selected columns, translated-row choices, and
\(\sigma\in S\),
\begin{equation}\label{eq:collision}
 |\det T|\le
  \exp\left\{-c\sum_a n_a^2+MH(E_{\rm hol}+\rho_H)\right\},
 \qquad \rho_H\longrightarrow0.
\end{equation}
One may take
\[
 \rho_H=
 \frac{\log M+c-\log(1-2^{-1/2})}{H}.
\]
\end{lemma}
\begin{proof}
On \(|z|\le R\), the column degree bound gives
\[
 |f_{a,P}(z)|
 \le\exp\left\{H\left(\frac R{w_0}
                        +\frac{\log(2R)}{w_*}\right)\right\}
 =:\mathcal B.
\]
On \(|t|=1/2\),
\[
 |j\lambda_\sigma+\log(1+t)|
 \le(K-1)\Gamma|x|+\log2<R/2.
\]
Expand \(f_{a,P}(z)=\sum_{d\ge0}b_{a,P,d}(z/R)^d\).
Cauchy's estimate gives \(|b_{a,P,d}|\le\mathcal B\).
The Taylor test of order \(r\) applied to \((z/R)^d\) has
absolute value at most \(2^r2^{-d}\).

If two rows use the same pair \((a,d)\), their coefficient vectors
\((b_{a,P,d})_P\) agree, and that determinant term is zero.
For a nonzero term the degrees in each \(a\)-group are distinct;
therefore
\[
 \sum_{\text{rows}}d\ge\sum_a\binom{n_a}{2}
                    =\frac{\sum_a n_a^2-M}{2}.
\]
Use half of the factor \(2^{-\sum d}\) for this lower bound and
sum its remaining half over all nonnegative degrees independently.
The latter sum is at most \((1-2^{-1/2})^{-M}\).
The determinant of each coefficient matrix is at most
\(M!\mathcal B^M\), and the sum of the Taylor-test orders is at most
\(MH/v_0\). Hence
\[
 |\det T|\le
 M!\mathcal B^M 2^{MH/v_0}
 \exp\{-c(\sum_a n_a^2-M)\}(1-2^{-1/2})^{-M}.
\]
Use \(M!\le M^M\) to obtain the asserted expression.
The series manipulations are justified by the same convergent
geometric majorant. Equation~\eqref{eq:row-count} shows
\(\log M/H\to0\). This proves the stated uniformity as well.
\end{proof}

The argument is the coefficient-collision mechanism of
\cite[Lemma 3.3]{OAIpi}. The change to real logarithmic centers
requires only the new radius bound; it does not remove the repeated
coefficient rows responsible for the saving.

\subsection{Two alternatives for the full determinant}

Fix \(0<A<1\) and \(0<\eta<1\), and let
\[
 D_A(H)=\#\{a\in\N^m:wa\le AH\},\qquad
 L_H=\frac{\eta^2M}{H D_A(H)}.
\]
The simplex counts imply
\begin{equation}\label{eq:collision-limit}
 D_A(H)\sim\frac{A^mH^m}{m!\prod_iw_i},\qquad
 L_H\longrightarrow
 L=\frac{\eta^2K\theta^m}{(m+1)v_0A^m}.
\end{equation}
If at least \(\eta M\) chosen rows have \(wa\le AH\),
Cauchy--Schwarz gives
\(\sum_a n_a^2\ge\eta^2M^2/D_A(H)=MH L_H\).
For a nonzero scalar, \(a\ge\beta\); thus the approximation
exponent in \eqref{eq:scalar} is nonpositive and may be discarded.
If fewer than \(\eta M\) rows have \(wa\le AH\), then
\[
 \sum_{\rm rows}w(a-\beta)>
 MH\bigl(A(1-\eta)-\bar b_H\bigr).
\]
In this alternative discard the nonpositive collision exponent.
The determinant expansion has at most \(Q_H^M\) terms.
Taking logarithms of the triangle inequality yields, at every
compatible embedding,
\begin{equation}\label{eq:good-bound}
 \frac{\log|\sigma(\Delta_H)|}{MH}
 \le E_{\rm tr}+E_{\rm hol}
     +\rho_H+\frac{\log Q_H}{H}
     +\max\{-cL_H,-\nu(A(1-\eta)-\bar b_H)\}.
\end{equation}
This upper bound applies to every selected nonzero full-row minor.
It is uniform in its selected columns; no analytic choice of a special
minor is made after the arithmetic choice.

\subsection{The other embeddings}

For an arbitrary embedding \(\sigma\), use
\eqref{eq:binomial-entry} directly. On \(|t|=1/2\),
\[
 |j\sigma(\gamma)p_i/q_i|\le KB_0,\qquad |G_i(t)|<1.
\]
The binomial coefficient costs at most \(2^{|\alpha|}\);
\(|\sigma(\xi)^{jh}|\le\Xi^{Kh}\);
\(|1+t|^h\le(3/2)^h\); and coefficient extraction costs \(2^\ell\).
It follows that every entry is at most \(\exp(HU)\), where the
convenient common bound used by the formal proof is
\begin{equation}\label{eq:direct-error}
 U=\frac{(\log\Xi+\log(3/2))K}{w_0}
       +\frac{\log(2(KB_0+2))}{w_*}
       +\frac{\log2}{v_0}\ge0.
\end{equation}
Thus
\begin{equation}\label{eq:other-bound}
 \frac{\log|\sigma(\Delta_H)|}{MH}\le U+\frac{\log M}{H}.
\end{equation}
We claim no collision saving at incompatible embeddings.

\begin{proposition}[All-embedding upper bound]\label{prop:analytic}
Let
\[
 E_{\rm an}=E_{\rm tr}+E_{\rm hol}+\kappa U.
\]
For the same minor used in Proposition~\ref{prop:arithmetic},
\begin{equation}\label{eq:analytic}
 \mathcal N_H\le E_{\rm an}+o(1)
       +\max\{-cL_H,-\nu(A(1-\eta)-\bar b_H)\}.
\end{equation}
The remainder tends to zero for the fixed data and is uniform over
the selected columns.
\end{proposition}
\begin{proof}
Sum \eqref{eq:good-bound} over \(S\), and
\eqref{eq:other-bound} over its complement, then divide by \(s\).
Since \(U+\log M/H\ge0\), replacing the complement size by \(d\)
only enlarges the upper bound.
One resulting remainder is
\[
 \rho_H+\frac{\log Q_H}{H}+\kappa\frac{\log M}{H},
\]
which tends to zero by the polynomial growth of \(M\) and \(Q_H\).
\end{proof}

\section{Parameter choices and the final contradiction}
\label{sec:parameters}

\begin{lemma}[The comparison gap]\label{lem:gap}
Let \(\kappa\ge1\) and \(\nu>2\kappa\). There are positive
rational numbers \(\theta,A,B,C\) and a real \(\eta\in(0,1)\)
such that
\begin{gather}
 0<\theta<A<B<1,\qquad C>1,\qquad
 C\theta<1,\quad CB<1,\quad C\theta/B>1,\quad B/A>1,
 \label{eq:parameter-ratios}\\
 g:=\nu\bigl(A(1-\eta)-\theta\bigr)-\kappa(1-\theta)>0.
 \label{eq:gap}
\end{gather}
\end{lemma}
\begin{proof}
Choose a rational \(b_0\) with
\[
 \frac12<b_0<1-\frac{\kappa}{\nu}.
\]
For a sufficiently small positive rational \(\delta\), put
\(\theta=1-\delta\), \(A=1-b_0\delta\). Then \(0<\theta<A<1\),
\[
 \nu(A-\theta)-\kappa(1-\theta)
     =\bigl(\nu(1-b_0)-\kappa\bigr)\delta>0,
\]
and
\(\theta-A^2=(2b_0-1)\delta-b_0^2\delta^2>0\).
Choose rationals
\[
 A/\theta<C<1/A,\qquad A<B<\min(1/C,C\theta).
\]
The intervals are nonempty because \(A^2<\theta\).
Also \(C\theta<\theta/A<1\), giving all the ratio inequalities.
Finally choose \(\eta>0\) small enough to preserve the positive gap.
\end{proof}

The constants in the preceding estimates can all be made small in the
required order. To see this without suppressing any number-field cost,
put \(E=E_{\rm ar}+E_{\rm an}\) and define
\[
 \begin{split}
 \mathcal C&=R_0+\kappa\log D_\xi+
                    \kappa(\log\Xi+\log(3/2)),\\
 \mathcal W&=\kappa\theta+\log4+\log(2K\Gamma)+\nu
       +\log(2R_0K)+\kappa\log(2(KB_0+2)).
 \end{split}
\]
All these quantities are nonnegative. The exact error identity is
\begin{equation}\label{eq:error-ledger}
 \begin{split}
 E={}&\frac{\nu}{F_0}
   +\frac{\kappa\Lambda F_0m+(\kappa+2)\log2}{v_0}
   +\mathcal C\frac K{w_0}\\
   &+\kappa(\Lambda+\log D_\gamma)\sum_{i=1}^m\frac1{w_i}
     +\frac{\mathcal W}{w_*}.
 \end{split}
\end{equation}
In particular, algebraic denominators and unconstrained conjugates
appear explicitly rather than being hidden in the \(o(1)\) in \(H\).

\begin{proposition}[Admissible finite choices]\label{prop:choices}
Assuming \eqref{eq:bad-approximations}, all fixed data can be selected
so that the interpolation theorem applies and
\begin{equation}\label{eq:margins}
 E<g,\qquad cL>\kappa+E.
\end{equation}
\end{proposition}
\begin{proof}
Choose \(\theta,A,B,C,\eta\) by Lemma~\ref{lem:gap}, and set
\(\epsilon_0=\min(g/2,1/2)>0\).
First choose \(F_0>2/\theta\) such that
\(\nu/F_0<\epsilon_0/3\).
For an integer \(m\ge1\), define
\begin{equation}\label{eq:dimension-parameters}
 K=\lfloor C^m\rfloor,\qquad w_0=B^{-m},\qquad
 v_0=2K\theta^m w_0.
\end{equation}
They satisfy
\[
 K(w_0/v_0)\theta^m=\tfrac12,\qquad
 K\theta^m\le(C\theta)^m<1.
\]
The second inequality allows the common-fiber branch as well as the
distinct-center branch.
Moreover
\[
 \frac K{w_0}\le(CB)^m\longrightarrow0,\qquad
 v_0\sim2(C\theta/B)^m\longrightarrow\infty,\qquad
 \frac m{v_0}\longrightarrow0.
\]
Substitution in \eqref{eq:collision-limit} gives the exact expression
\begin{equation}\label{eq:L-dimension}
 L=\frac{\eta^2}{2(m+1)}(B/A)^m\longrightarrow\infty.
\end{equation}
Thus one can fix a finite \(m\) large enough that
\[
 \frac{\kappa\Lambda F_0m+(\kappa+2)\log2}{v_0}
       +\mathcal C\frac K{w_0}<\epsilon_0/3,
 \qquad cL>\kappa+1.
\]
This fixes \(m,K,w_0,v_0\), and hence \(\mathcal W\).

Now choose \(\sigma_0\) and the weight thresholds for
Theorem~\ref{thm:interpolation}, imposing also its common-fiber
volume inequality. Choose the finitely many approximants
\(p_i/q_i\) successively from \eqref{eq:bad-approximations}.
Their weights can satisfy every separation threshold and can all
exceed a common bound large enough that
\[
 \kappa(\Lambda+\log D_\gamma)\sum_i\frac1{w_i}
               +\frac{\mathcal W}{w_*}<\epsilon_0/3.
\]
Increase the denominator thresholds as needed to ensure \(q_i\ge2\)
and \(p_i\ne0\).
This step is valid because the interpolation weight thresholds are
independent of the centers, which themselves involve \(p_i/q_i\).

Equation~\eqref{eq:error-ledger} now gives \(E<\epsilon_0<g\);
also \(cL>\kappa+1>\kappa+E\).
Every choice is finite. Only after fixing this entire list do we let
the sufficiently divisible parameter \(H\) increase.
\end{proof}

\begin{proof}[Proof of Theorem~\ref{thm:compatible}]
Normalize \((\gamma,\xi)\) by Lemma~\ref{lem:torsion}.
Fix \(\nu>2\kappa\) and suppose \eqref{eq:bad-approximations}.
Make the choices of Proposition~\ref{prop:choices}.
For arbitrarily large divisible \(H\), interpolation supplies a
nonzero selected determinant. Propositions~\ref{prop:arithmetic} and
\ref{prop:analytic} then give
\[
 -\kappa(1-\bar b_H)-E_{\rm ar}
 \le E_{\rm an}+o(1)
       +\max\{-cL_H,-\nu(A(1-\eta)-\bar b_H)\}.
\]
Both entries of this maximum are incompatible with the lower bound
for sufficiently large \(H\).
For the collision entry, \(0\le\bar b_H\) implies
\(\kappa(1-\bar b_H)\le\kappa\), whereas
\(cL_H\to cL>\kappa+E\).
For the other entry, since \(\nu>\kappa\) and
\(\bar b_H\le\theta\),
\[
 \nu\bigl(A(1-\eta)-\bar b_H\bigr)
     -\kappa(1-\bar b_H)
 \ge \nu\bigl(A(1-\eta)-\theta\bigr)-\kappa(1-\theta)
 =g>E.
\]
The strict margins absorb the common remainder \(o(1)\).
This contradiction excludes \eqref{eq:bad-approximations}.
It proves \(\ELB(x,2\kappa)\), including the bound for unreduced
rational representations. Lemma~\ref{lem:endpoint} proves
irrationality and the exponent bounds.
The rational-logarithm theorem and algebraic-base corollary follow
from the specializations already given.
\end{proof}

\begin{remark}[What the exponent threshold means]
The condition \(\nu>2\kappa\) is used to choose
\(1/2<b_0<1-\kappa/\nu\).
It supplies the two simultaneous requirements \(A^2<\theta\) and a
positive determinant gap. This is the precise source of \(2d/s\).
At the boundary \(\nu=2\kappa\), this argument gives neither that
interval nor a strict gap.
For rational logarithms, exponent two means the eventual inequalities
for every \(2+\varepsilon\); it does not assert an eventual lower
bound with the exact exponent \(2\).
\end{remark}

\section{Special cases, comparison with the \texorpdfstring{\(\pi\)}{pi}
proof, and limitations}\label{sec:consequences}

\subsection{Rational powers and rational affine changes}

\begin{lemma}\label{lem:affine}
If \(b\ge2\), \(\ELB(x,b)\), and \(a,c\in\Q\) with \(a\ne0\),
then \(\ELB(ax+c,b)\).
\end{lemma}
\begin{proof}
For a fixed \(\nu>b\), choose \(b<\nu'<\nu\).
Write \(a=A/B\), \(c=C/D\), with \(B,D>0\).
For any \(p/q\), the rational number \((p/q-c)/a\) has an integer
representation with denominator \(|A|Dq\).
Consequently, for sufficiently large \(q\),
\[
 |ax+c-p/q|\ge |a|(|A|Dq)^{-\nu'}.
\]
The fixed positive factor
\(|a|(|A|D)^{-\nu'}\) is eventually at least
\(q^{-(\nu-\nu')}\), proving the assertion.
\end{proof}

\begin{corollary}\label{cor:rational-powers}
Suppose \(\alpha>0\), \(n\in\N_{>0}\), and
\(\alpha^n=r\in\Q_{>0}\setminus\{1\}\).
Then
\[
 \mu(\log\alpha)=2.
\]
In particular,
\(\mu(\log\sqrt2)=2\).
\end{corollary}
\begin{proof}
The identity \(\log\alpha=(1/n)\log r\), together with
Theorem~\ref{thm:rational}, Lemma~\ref{lem:affine}, and
Lemma~\ref{lem:endpoint}, proves the result.
\end{proof}

This sharp conclusion is stronger than the degree bound in this
particular family. The proof uses rational scaling, not invariance
under arbitrary algebraic scaling.

\subsection{Arctangents of rational numbers}

\begin{corollary}\label{cor:arctangent}
For every \(r\in\Q\setminus\{0\}\),
\[
 \mu(\arctan r)=2.
\]
\end{corollary}
\begin{proof}
Use \(F=\Q(i)\), \(\gamma=2i\),
\(\xi=(1+ir)/(1-ir)\), and \(x=\arctan r\ne0\).
The identity
\[
 e^{2i\arctan r}=\frac{1+ir}{1-ir}
\]
establishes compatibility for the given embedding. Its complex
conjugate establishes compatibility for the other embedding.
Thus \(s=d=2\).
For \(r=1\), \(\xi=i\) is a root of unity; the root-of-unity
branch of the theorem is genuinely needed here.
\end{proof}

\subsection{Quadratic-normalized logarithms}

\begin{corollary}\label{cor:quadratic}
Let \(a,b,D\in\Q\), \(D>0\), and suppose
\[
 a^2-b^2D=1,\qquad \alpha=a+b\sqrt D>1.
\]
Then
\[
 \mu\left(\frac{\log(a+b\sqrt D)}{\sqrt D}\right)=2.
\]
In particular,
\begin{equation}\label{eq:quadratic-example}
 \mu\left(\frac{\log(3+2\sqrt2)}{\sqrt2}\right)=2.
\end{equation}
\end{corollary}
\begin{proof}
Take \(F=\Q(\sqrt D)\), \(\gamma=\sqrt D\),
\(\xi=a+b\sqrt D\), and \(x=\log\alpha/\sqrt D\ne0\).
Every embedding sends \(\sqrt D\) to either \(\sqrt D\) or
\(-\sqrt D\). The first choice gives
\(e^{\sqrt D\,x}=\alpha\).
The norm identity gives \(a-b\sqrt D=\alpha^{-1}\), so the
second choice gives
\[
 e^{-\sqrt D\,x}=\alpha^{-1}=a-b\sqrt D.
\]
Every embedding is compatible, whether the field has degree one or two.
\end{proof}

The normalization in this corollary is essential to the asserted
specialization. It does not imply
\(\mu(\log(3+2\sqrt2))=2\) by removing the factor \(\sqrt2\).
The general algebraic-base result gives the separate bound
\(2\le\mu(\log(3+2\sqrt2))\le4\).

\subsection{How the argument differs from the original
\texorpdfstring{\(\pi\)}{pi} proof}

The relation to \cite{OAIpi} is at the level of proved machinery,
not a substitution into its final theorem.
The following comparison fixes the main changes.

{\small
\begin{longtable}{@{}>{\raggedright\arraybackslash}p{.22\textwidth}
>{\raggedright\arraybackslash}p{.34\textwidth}
>{\raggedright\arraybackslash}p{.36\textwidth}@{}}
\toprule
Feature & Original \(\pi\) construction & Logarithmic extension\\
\midrule
\endfirsthead
\toprule
Feature & Original \(\pi\) construction & Logarithmic extension\\
\midrule
\endhead
\bottomrule
\endfoot
Exponential identity
& \(e^{2i\pi}=1\)
& \(e^{\log r}=r\), or \(e^{\sigma(\gamma)x}=\sigma(\xi)\)\\[5pt]
Multiplicative centers
& Common fiber \(Y=1\)
& Distinct powers \(r^j\) or \(\xi^j\); common fiber after
root-of-unity reduction\\[5pt]
Local coordinates
& \(t=Y-1\)
& \(t=Y/y_j-1\); the same logarithmic frame\\[5pt]
Arithmetic
& Gaussian-integer norm after denominator clearing
& Rational integers for rational logarithms; an algebraic integer
and its full field norm in general\\[5pt]
Analytic saving
& Repeated \((a,d)\) coefficient rows
& The same repetition argument at every compatible embedding;
direct growth bounds at the others\\[5pt]
Final coefficient
& \(1\) in the normalized arithmetic comparison
& \(d/s\), giving threshold \(2d/s\)\\
\end{longtable}
}

For \(\pi\), the compatible data may be viewed as
\(F=\Q(i)\), \(\gamma=2i\), \(\xi=1\), \(x=\pi\):
both embeddings are compatible. For \(\log r\), the field is
\(\Q\), the slope is \(1\), and the real centers require the
distinct-fiber argument. The analytic translation is preserved
exactly, while the radius constants change.
The proof of the new interpolation theorem, the identification with
\eqref{eq:matrix}, the field-denominator estimates, and the estimates
at all embeddings are indispensable additions.

\subsection{Why no theorem for all positive real bases is possible}

\begin{proposition}\label{prop:counterexamples}
The map \(\alpha\mapsto\log\alpha\) on positive real numbers is
surjective onto \(\R\). There is a positive real \(\alpha\ne1\)
with rational logarithm, and there is such an \(\alpha\) whose
irrational logarithm has no finite irrationality exponent.
\end{proposition}
\begin{proof}
For every \(y\in\R\), \(\alpha=e^y>0\) satisfies
\(\log\alpha=y\). Taking \(y=1\) gives the rational example.
For the second, let
\[
 L=\sum_{n=0}^\infty10^{-n!},\qquad \alpha=e^L.
\]
The standard Liouville argument applies: truncation through \(N\)
has denominator \(q_N=10^{N!}\), while its positive tail is at most
\(2q_N^{-(N+1)}\). These approximations prove that \(L\) is an
irrational Liouville number and admits no eventual lower bound of
any finite exponent. Thus \(\log\alpha=L\) has the claimed property.
\end{proof}

These examples concern arbitrary real inputs. They are not
counterexamples to the rational- or algebraic-base theorems.
The formal statements express the second conclusion as failure of
\(\ELB(\log\alpha,b)\) for every finite real \(b\), avoiding any
ambiguous convention for a real-valued supremum of an unbounded set.

\section{Formal verification and reproducibility}\label{sec:verification}

\subsection{What the formal endpoints assert}

The Lean definition \path{LogarithmExtension.EventualLowerBound}
is exactly the predicate \(\ELB\) in
Section~\ref{sec:preliminaries}; it quantifies over every integer
numerator and every sufficiently large natural denominator.
The main declarations include
\begin{Verbatim}[fontsize=\small]
logThree_eventual_lower_bound
logThree_irrationalityExponent_eq_two
rationalLog_eventual_lower_bound
rationalLog_irrationalityExponent_eq_two
algebraicLog_irrationalityExponent_bounds
CompatibleLog.compatible_embedding_eventual_lower_bound
CompatibleLog.compatible_embedding_irrationalityExponent_bounds
CompatibleLog.all_embeddings_irrationalityExponent_eq_two
\end{Verbatim}
All are in the namespace \texttt{LogarithmExtension}.
The unbundled compatible-embedding theorem takes only the number field,
\(\gamma,\xi,x\), their stated nonzero conditions, a nonempty finite
set of embeddings, and the exponential identities
\eqref{eq:compatibility}. It has no interpolation, nonvanishing,
analytic-bound, or irrationality premise.
The bundled data structure \path{CompatibleLog.Base} contains these
same mathematical data, not a hidden proof of the conclusion.

The names \path{LogThreeTarget}, \path{RationalLogarithmTarget},
and \path{AlgebraicLogarithmTarget} are proposition definitions.
They are not axioms: the corresponding \texttt{target} declarations
prove them.

The exponent definition reused from \texttt{OAI.PiExponent} is a
real supremum of the positive exponents admitting infinitely many
good reduced rational approximants. For every finite-bound endpoint,
the code first proves that this set is bounded above, and that it
contains \(2\). The resulting supremum therefore agrees with the
usual extended-real definition in this paper.
The Liouville counterexample is expressed as absence of any finite
eventual bound, rather than as a possibly misleading real supremum.

\subsection{Correspondence of arguments and source files}

All extension paths below are relative to
\texttt{LogarithmExtension/}. The earlier dedicated
\texttt{LogThree}, \texttt{RationalLog}, and \texttt{AlgebraicLog}
pipelines remain checked; the \texttt{CompatibleLog} pipeline supplies
the unified theorem proved here.

{\small
\let\path\leanfile
\begin{longtable}{@{}>{\raggedright\arraybackslash}p{.37\textwidth}
                     >{\raggedright\arraybackslash}p{.59\textwidth}@{}}
\toprule
Mathematical component & Principal extension modules\\
\midrule
\endfirsthead
\toprule
Mathematical component & Principal extension modules\\
\midrule
\endhead
\bottomrule
\endfoot
Actual target and finite exponent semantics
& \path{Target.lean}, \path{Endpoint.lean},
\path{GeneralExponent.lean}\\[4pt]
Normalized logarithmic coordinates and rigidity
& \path{NormalizedLogJet.lean}, \path{CoordinateScaling.lean},
\path{CurveRigidity.lean}\\[4pt]
Proposition~\ref{prop:curve}, distinct-center case
& \path{CurveContactFamily.lean},
\path{CurveDerivativeVanishing.lean}, \path{CurveInequality.lean},
\path{CurveGeometry.lean}\\[4pt]
Common-fiber branch
& \path{CommonFiberCurve.lean}, \path{GeometryData.lean}\\[4pt]
Projective degree, local ideals, and exceptional degree
& \path{CurveModelDegrees.lean}, \path{NormalizedJetIdeal.lean},
\path{CompactNormalizedJetIdeal.lean}, \path{LocalContactLength.lean},
\path{ExceptionalDegree.lean}\\[4pt]
Ampleness and global jet generation
& \path{BlowupGeometry.lean}, \path{BlowupMargin.lean},
\path{BlowupAmple.lean}, \path{JetAmpleness.lean},
\path{GlobalJetPackets.lean}, \path{JetSurjectivity.lean}\\[4pt]
Identification with the actual matrix
& \path{BoundedJetMatrix.lean}, \path{MatrixIdentification.lean},
\path{CompatibleLogGeometry.lean}\\[4pt]
Equation~\eqref{eq:row-expansion} and analytic centers
& \path{AnalyticTranslation.lean},
\path{CompatibleLogAnalyticCenters.lean}\\[4pt]
Arithmetic over \(F\), integral norm, and averaging
& \path{NumberFieldMatrix.lean}, \path{NumberFieldNorm.lean},
\path{NumberFieldDeterminant.lean},
\path{NumberFieldNormAveraging.lean},
\path{CompatibleLogMatrixArithmetic.lean}\\[4pt]
Actual summands and their finite sum
& \path{CompatibleLogLiteralAnalyticSummand.lean},
\path{CompatibleLogAnalyticAggregate.lean}\\[4pt]
Incompatible embeddings and full norm bound
& \path{ConjugateAnalyticBound.lean},
\path{CompatibleLogAllEmbeddings.lean}\\[4pt]
Choices and determinant contradiction
& \path{CompatibleLogSuccessiveApproximations.lean},
\path{CompatibleLogAdmissibleParameters.lean},
\path{CompatibleLogDeterminantContradiction.lean}\\[4pt]
Root-of-unity reduction and final theorem
& \path{CompatibleLogBase.lean}, \path{CompatibleLogNonTorsion.lean},
\path{CompatibleLogMain.lean}\\[4pt]
Theorem~\ref{thm:rational}, including the literal \(\log3\)
& \path{LogThreeMain.lean}, \path{RationalLogMain.lean}\\[4pt]
Corollaries~\ref{cor:algebraic}--\ref{cor:quadratic}
& \path{AlgebraicLogMain.lean}, \path{RationalAffine.lean},
\path{RationalPowerLog.lean}, \path{RationalArctangent.lean},
\path{QuadraticNormalizedLog.lean}\\[4pt]
Proposition~\ref{prop:counterexamples}
& \path{LogarithmCounterexamples.lean}\\
\end{longtable}
}

General results reused from the pinned upstream development include
the persistent-component and weighted multiplicity machinery,
projective and local algebra, numerical ampleness and vanishing,
simplex counts, logarithmic row algebra, collision estimates, and
the approximation endpoint lemmas. Theorem \(\mu(\pi)=2\) itself
is not the hypothesis from which the logarithm result is deduced.
The new target-independent geometry data remove the
\(\pi\)-specific approximation hypotheses from the geometric interface.

\subsection{Axiom audit and kernel replay}

The build and transitive axiom audit use the following commands:
\begin{Verbatim}[fontsize=\small]
lake build -q LogarithmExtension.Verification
lake env lean --trust=0 -Ddebug.skipKernelTC=false \
  LogarithmExtension/Verification.lean
\end{Verbatim}
The audit checked \(1326\) extension theorem declarations.
Every transitive axiom was one of
\[
 \texttt{propext},\qquad\texttt{Classical.choice},\qquad
 \texttt{Quot.sound}.
\]
It also audited the reference \(\pi\) endpoint separately.
A source scan found no proof admissions, new axiom declarations,
unsafe or partial declarations, or \texttt{native\_decide} in the
extension. A defined target proposition is not counted as a proof.

The full combined replay was also rerun successfully:
\begin{Verbatim}[fontsize=\small]
lake env lean --run VerifyGeneralLog.lean
\end{Verbatim}
Its observed results were
\begin{Verbatim}[fontsize=\small]
Replaying 118179 endpoint dependencies in an empty kernel environment.
PASS: all 33 general-logarithm and counterexample endpoints.
PASS: their complete dependency closure rechecked
      in an empty kernel environment.
PASS: only propext, Classical.choice, and Quot.sound allowed as axioms.
\end{Verbatim}
The selected roots cover \(\log3\), rational and rational-power
logarithms, algebraic logarithms, both branches of the compatible
theorem, the arctangent and quadratic examples, and the real-input
counterexamples.
The checker collects dependencies of both types and values, follows
mutual declaration blocks, rejects unsafe or partial declarations
and nonstandard axioms, calls \texttt{Lean.Replay} in an environment
created by \texttt{mkEmptyEnvironment}, and checks that the selected
roots occur in the replayed environment.
This replay uses Lean's kernel in a fresh environment, not a second,
independently implemented proof checker.

\subsection{Pinned sources and reproduction}

The source package fixes the following revisions:
\begin{center}
\small
\begin{tabular}{@{}ll@{}}
\toprule
Component & Version or commit\\
\midrule
Lean & \texttt{4.34.1}\\
Lean commit & \texttt{5045d0056413266e57c625dcd7c365b10e377c52}\\
Mathlib & \texttt{d13f23b723b8a846827a245b89c10fc7d3f11612}\\
OpenAI mathematics & \texttt{adc7f1241b42e322a6451854ab7e4b4c146bf78a}\\
\bottomrule
\end{tabular}
\end{center}
The upstream tracked source was not modified.
The ancillary package contains the extension, selected upstream
source trees, \texttt{lean-toolchain}, \texttt{lakefile.toml},
the pinned \texttt{lake-manifest.json}, and the replay programs.
A SHA-256 inventory identifies the packaged files.
Only local source modules required by the logarithm formalization and
its verification targets are included.

After installing the specified Lean toolchain, enter the ancillary
Lean project directory and run
\begin{Verbatim}[fontsize=\small]
lake exe cache get
lake build LogarithmExtension.Verification
lake env lean --run VerifyGeneralLog.lean
\end{Verbatim}
Finish the build before starting the replay, and do not rebuild
imported modules while it is running.
The cache command obtains the Mathlib build cache; compiling its
sources is an alternative.
External toolchain and package downloads, adequate memory, and disk
space are required. The distributed archive contains sources rather
than a preinstalled Lean environment.

\subsection*{Acknowledgments}

The determinant method and much of its formal infrastructure are due
to the OpenAI preprint \cite{OAIpi}. Their reuse and adaptation are
credited throughout, and the upstream license notices are retained
in the source package.
The Lean and Mathlib communities supplied the proof assistant and
general mathematical libraries.
This work was developed with substantial assistance from OpenAI's
Codex, including exploration, implementation of the formal proofs,
and manuscript preparation.

\bibliographystyle{plain}
{\raggedright
\bibliography{references}
}
\end{document}
